% Motivating example (Fig. 1): three semantic subfigures.
\begin{subfigure}[t]{0.32\linewidth}
  \centering
  \begin{tikzpicture}[x=0.72cm,y=0.72cm]
    \node[mannode] (X1) at (0,1.15) {$X_1$};
    \node[mannode] (X2) at (1.5,1.15) {$X_2$};
    \node[tgtnode] (Y) at (0.75,-0.35) {$Y$};
    \draw[diredge] (X1) -- (Y);
    \draw[diredge] (X2) -- (Y);
    \draw[bicross] (X1) to[bend left=28] (X2);
  \end{tikzpicture}
  \caption{True graph}
  \label{fig:motivating-true}
\end{subfigure}\hfill
\begin{subfigure}[t]{0.32\linewidth}
  \centering
  \begin{tikzpicture}[x=0.72cm,y=0.72cm]
    \node[mannode] (X1) at (0,1.15) {$X_1$};
    \node[mannode] (X2) at (1.5,1.15) {$X_2$};
    \node[tgtnode] (Y) at (0.75,-0.35) {$Y$};
    \draw[diredge,black!25] (X1) -- (Y)
      node[midway,sloped] {\textcolor{red!70!black}{\footnotesize$\boldsymbol{\times}$}};
    \draw[diredge] (X2) -- (Y);
    \draw[bicross] (X1) to[bend left=28] (X2);
  \end{tikzpicture}
  \caption{Assumed graph}
  \label{fig:motivating-assumed}
\end{subfigure}\hfill
\begin{subfigure}[t]{0.32\linewidth}
  \centering
  \begin{tikzpicture}[x=0.72cm,y=0.72cm]
    \node[mannode] (X1) at (0,1.15) {$X_1$};
    \node[mannode] (X2) at (1.5,1.15) {$X_2$};
    \node[tgtnode] (Y) at (0.75,-0.35) {$Y$};
    \begin{scope}[on background layer]
      \node[clbox,fit=(X1)(X2),label={[font=\scriptsize]above:$C_1$}] (C1) {};
    \end{scope}
    \draw[diredge] (C1.south) -- (Y);
  \end{tikzpicture}
  \caption{Common quotient}
  \label{fig:motivating-quotient}
\end{subfigure}
